% Minimal retained preamble for lane-p2 cards taken from
% sources/logic-fl02-conservation-of-ramsey-s-theorem-f/pinned-original-source.tex
% (Conservation of Ramsey's theorem for pairs, arXiv:2402.11616v3).
%
% Copied verbatim from the author's own file:
%   line 2      : \documentclass{article}
%   lines 16-25 : the author's theorem declarations
%   lines 46-88 : the author's macro definitions for the reverse-mathematics principles
%   line 91     : \card
%   lines 104-107: the natural sums and products, which rely on mathabx (source line 103)
%   line 116    : \qt
%
% Omitted author preamble lines and why (nothing was renamed or redefined):
%   line 3 (graphicx), 7 (url), 9 (cancel): no figure, url or cancel use in the retained spans;
%   line 6 (cleveref): the assembler adds its own cleveref machinery, a second load clashes;
%   line 8 (\usepackage[utf8x]{inputenc}): inputenc with the utf8x option fights the forced
%       xelatex engine;
%   lines 28-37 (the rep@theorem machinery and \newreptheorem{maintheorem}): the retained spans
%       use no rep-theorem environment;
%   lines 40-43 (xcolor and the three author comment macros): not used by the retained spans;
%   lines 93-99 (the LS2 font encoding block for \smashtimes based on stix): the host TeX Live has
%       no stix directory and the author himself never uses the result -- line 105 maps \ntimes to
%       \dottimes, and line 110 shows \smashtimes already commented out in the source.
% mathabx is retained because \dotplus, \dottimes and \dotdiv come from it.
\documentclass{article}
\usepackage{amsmath,amssymb,amsthm}
\usepackage{mathabx}

\newtheorem{theorem}{Theorem}
\numberwithin{theorem}{section}
\newtheorem{maintheorem}[theorem]{Main Theorem}
\newtheorem{lemma}[theorem]{Lemma}
\newtheorem{question}[theorem]{Question}
\newtheorem{proposition}[theorem]{Proposition}
\newtheorem{statement}[theorem]{Statement}
\newtheorem{remark}[theorem]{Remark}
\newtheorem{definition}[theorem]{Definition}
\newtheorem{corollary}[theorem]{Corollary}

\newcommand{\RCA}{\mathsf{RCA}}
\newcommand{\WKL}{\mathsf{WKL}}
\newcommand{\SRT}{\mathsf{SRT}}
\newcommand{\EM}{\mathsf{EM}}
\newcommand{\SEM}{\mathsf{SEM}}
\newcommand{\ADS}{\mathsf{ADS}}
\newcommand{\SCAC}{\mathsf{SCAC}}
\newcommand{\SADS}{\mathsf{SADS}}
\newcommand{\COH}{\mathsf{COH}}
\newcommand{\ACA}{\mathsf{ACA}}
\newcommand{\qvdash}{\operatorname{{?}{\vdash}}}
\newcommand{\nqvdash}{\operatorname{{?}{\nvdash}}}
\newcommand{\BSig}{\mathsf{B}\Sigma^0}
\newcommand{\ISig}{\mathsf{I}\Sigma^0}
\newcommand{\LSig}{\mathsf{L}\Sigma^0}
\newcommand{\BPi}{\mathsf{B}\Pi^0}
\newcommand{\IPi}{\mathsf{I}\Pi^0}
\newcommand{\LPi}{\mathsf{L}\Pi^0}
\newcommand{\IDel}{\mathsf{I}\Delta^0}
\newcommand{\BCDe}{\mathsf{BC}\Delta^0}
\newcommand{\WF}{\mathsf{WF}}
\newcommand{\RT}{\mathsf{RT}}
\newcommand{\BME}{\mathsf{BME}}
\newcommand{\TT}{\mathsf{TT}}
\renewcommand{\O}{\mathcal{O}}
\renewcommand{\L}{\mathcal{L}}
\newcommand{\C}{\mathcal{C}}
\newcommand{\R}{\mathcal{R}}
\renewcommand{\S}{\mathcal{S}}
\newcommand{\T}{\mathcal{T}}
\renewcommand{\P}{\mathcal{P}}
\newcommand{\M}{\mathcal{M}}
\newcommand{\Nc}{\mathcal{N}}
\renewcommand{\P}{\mathcal{P}}
\newcommand{\Q}{\mathcal{Q}}
\newcommand{\NN}{\mathbb{N}}
\newcommand{\PP}{\mathbb{P}}
\newcommand{\uh}[0]{{\upharpoonright}}
\newcommand{\card}{\operatorname{card}}
\newcommand{\nplus}{\dotplus}
\newcommand{\ntimes}{\dottimes}
\newcommand{\nminus}{\dotdiv}
\newcommand{\nsum}{\dot{\sum}}
\def\qt#1{``#1''}
\title{A condition can force every functional below k to fail being a long alpha-decreasing sequence}
\author{Quentin Le Houérou, Ludovic Levy Patey, Keita Yokoyama}
\date{arxiv:2402.11616v3; CC BY 4.0}
\begin{document}
\maketitle
\noindent\textit{Editorial note.} The counted unit is the article's own Lemma lem:coh-wf-forcing-wf at source lines 522-524 together with the proof environment that begins on the next line (525-543); the pairing is the ordinary adjacent one.\par\medskip
\noindent\textit{Editorial note.} The retained context at source lines 480-491 is the author's own account of conditions, of compatibility of a string with a condition and of the two forcing clauses for the existential and universal jump formulas; the retained context at source line 505 is his own explanation of the relation forcing a decided initial segment of the jump.\par\medskip
\noindent\textit{Editorial note.} The retained preamble reproduces verbatim the author's document class at source line 2, his theorem declarations at source lines 16-25, his macro definitions at source lines 46-88 and 91, his natural-operation macros at source lines 104-107 and his quotation macro at source line 116. The LS2 font-encoding block at source lines 93-99, which declares \textbackslash{}smashtimes through the stix fonts, is not retained because the host TeX Live has no stix directory and because the author never uses that symbol: source line 105 defines \textbackslash{}ntimes as \textbackslash{}dottimes and source line 110 shows the alternative already commented out. mathabx (source line 103) is retained because \textbackslash{}dotplus and \textbackslash{}dottimes come from it. The cleveref load at source line 6 and the inputenc load with the utf8x option at source line 8 are not retained, for the assembler's own cleveref machinery and the forced xelatex engine respectively. Nothing was renamed, redefined or added.\par\medskip
\noindent\textit{Editorial note.} No retained line contains a citation or a cross-reference, so the material carries no bibliography and needs no navigation pointer.\par\medskip
\noindent\textit{Editorial note.} The article's conservation theorems, its lemma lem:shore-blocking-wf, its lemma lem:wf-product, its lemma lem:coh-wf-forcing-jump and its proposition prop:coh-wf are not counted here and are not used as premises.\par\medskip
\noindent\textit{Allowed standard prerequisites (editorial).}\begin{itemize}
\item Conditions of the Cohen forcing used for COH. A condition is a pair (sigma, rho) with sigma a finite binary string and rho an M-finite approximation of the sequence of sets to be made cohesive, and it may be read as a Mathias condition; a string tau is compatible with (sigma, rho) when sigma is an initial segment of tau and the difference tau minus sigma is contained in the members of the approximation; the article's own account is retained at source lines 480-483.. Fixes the two objects of the counted lemma, namely a condition (sigma,rho) and the compatibility relation over which the recursive construction runs.
\item The forcing relation for the existential and universal jump formulas. A condition forces the existential statement that some value of the functional on the generic is defined when a witness and a bound smaller than the length of sigma exist, and it forces the universal statement that all values diverge when every string compatible with the condition fails to converge anywhere; the article's own definition is retained at source lines 485-491, and its explanation of forcing a decided initial segment of the jump at source line 505.. The counted lemma concludes a forcing statement about the generics G and Y, so the exact meaning of the forcing relation and of compatibility must be available; the proof uses the two clauses in its contradiction argument and its final transfer from conditions to generics.
\item Theories Sigma\textasciicircum{}0\_2 collection, Ramsey's theorem for singletons and well-foundedness. The proof uses that a model M satisfying the collection scheme for Sigma\textasciicircum{}0\_2 formulas satisfies the Ramsey theorem for singletons with k colours for every k, and that M satisfying WF(alpha) means that M has no M-infinite sequence of ordinals below alpha that decreases at every step; the article states both notations and uses them throughout.. The proof obtains an M-infinite homogeneous set for the colours given by the indices below k from the collection scheme, and the contradiction is the existence of an M-infinite alpha-decreasing sequence, which WF(alpha) forbids.
\end{itemize}\medskip
One can think of a condition $(\sigma, \rho)$ as a Mathias condition $(\sigma, \vec{R}_\rho)$. Requiring that $\rho \prec Q$ ensures that $\vec{R}_\rho$ is $M$-infinite. When increasing the length of~$\rho$, we make progress in cohesiveness.
Note that being a condition is a $\Delta^0_1(Q)$ predicate, as it simply requires checking that $\sigma \prec Q$, while the extension relation is $\Delta^0_1(Y)$. In particular, both are $\Delta^0_1(Q \oplus Y)$, so $\Delta^0_1(P)$.

A sequence $\tau$ is said to be \emph{compatible} with a condition $(\sigma,\rho)$ if $\sigma \prec \tau$ and  $\tau - \sigma \subseteq \vec{R}_{\rho}$.

\begin{definition}
Let~$(\sigma, \rho)$ be a condition, and $\Phi_e$ be a Turing functional.
\begin{enumerate}
    \item[1.] $(\sigma, \rho) \Vdash \exists x \Phi^{G \oplus Y}_e(x)\downarrow$ if there is some~$x, t < |\sigma|$ such that $\Phi^{\sigma \oplus Y}_e(x)[t]\downarrow$
    \item[2.] $(\sigma, \rho) \Vdash \forall x \Phi^{G \oplus Y}_e(x)\uparrow$ if for every~$\tau$ compatible with~$(\sigma, \rho)$ and every~$x \in M$, $\Phi^{\tau \oplus Y}_e(x)\uparrow$.
\end{enumerate}
\end{definition}

We will need the following two technical lemmas, which are Shore blocking arguments for satisfying the two kind of requirements, namely, deciding the jump and preserving $\WF(\alpha)$. In what follows, $(\tau, \rho) \Vdash (G \oplus Y)' \uh_k = \sigma'$ means that for every~$e < |\sigma'|$, if $\sigma'(e) = 1$ then $(\tau, \rho) \Vdash \Phi^{G \oplus Y}_e(e)\downarrow$ and if $\sigma'(e) = 0$, then $(\tau, \rho) \Vdash \Phi^{G \oplus Y}_e(e)\uparrow$. Thus, this relation is $\Pi^0_1(Y)$.

\begin{lemma}\label[lemma]{lem:coh-wf-forcing-wf}
Let $(\sigma, \rho)$ be a condition. For every~$k \in M$, there exists some~$n \in M$ an extension $(\tau, \rho) \leq (\sigma, \rho)$ forcing~$\Phi_e^{G \oplus Y}$ not to be an $\alpha$-decreasing sequence of length more than~$n$ for every~$e < k$.
\end{lemma}

\begin{proof}
By twisting the Turing functionals, we can assume that for every~$e, n \in M$, if $\Phi^\sigma_e(n)\downarrow$, then
\begin{itemize}
    \item $\Phi^\sigma_e(m)\downarrow$ for every~$m < n$
    \item $\Phi^\sigma_e(0), \Phi^\sigma_e(1), \dots, \Phi^\sigma_e(m)$ is a strictly $\alpha$-decreasing sequence
\end{itemize}

Suppose the lemma does not hold. Then in particular, for every~$M$-finite string~$\tau$ compatible with $(\sigma, \rho)$ and every~$n \in M$, there is some~$\tau_1 \succeq \tau$ compatible with~$(\sigma, \rho)$ and some~$e < k$ such that $\Phi^{\tau_1 \oplus Y}(n)\downarrow$.

Build $Y$-recursively an infinite sequence $\tau_0 \preceq \tau_1 \preceq \tau_2 \preceq \dots$ such that for every~$n$,
\begin{itemize}
    \item $\tau_n$ is compatible with $(\sigma, \rho)$
    \item $\Phi^{\tau_n \oplus Y}_e(n)\downarrow$ for some~$e < k$ (which depends on~$n$)
\end{itemize}
Since~$\M \models \BSig_2$, $\M \models \forall k\RT^1_k$,
so there is an $M$-infinite set~$Z \in \M$ and some~$e < k$ such that for every~$n \in Z$, $\Phi^{\tau_n \oplus Y}_e(n)\downarrow$.

Then, letting $f(n) = \Phi^{\tau_n \oplus Y}_e(n)$, we have an $M$-infinite $\alpha$-decreasing sequence, contradicting $\M \models \WF(\alpha)$.
\end{proof}

\end{document}
